% ============================================================
% COMPÉTENCES — VARIANTE UNE COLONNE
% ============================================================
% Mêmes compétences que sections/sidebar/skills.tex, mais à plat : une ligne
% par catégorie plutôt qu'une liste à puces, la pleine largeur rendant les
% listes verticales inutilement longues.
%
% Règle tenue ici : chaque techno listée est adossée à quelque chose de
% vérifiable dans le corps du CV. Go n'apparaît que dans les projets
% personnels, d'où la mention explicite ; Node, qui n'était rattaché à aucune
% expérience, a été retiré. Mercure est un transport temps réel, il est classé
% en Outils et non en Données.
%
% Les langues sont rattachées ici plutôt qu'à une section dédiée en fin de CV :
% deux lignes de contenu ne justifiaient pas un titre avec filet, et le tableau
% se lit comme un inventaire de ce que la personne sait faire — les langues y
% ont leur place.
\section*{\faIcon{tools} Compétences}

\cvskillline{Back-end}{PHP / Symfony / API Platform, Go \textit{(projets personnels)}}
\cvskillline{Front-end}{React, Vue, React Native, Twig}
\cvskillline{Données}{SQL, SQLite, Doctrine, API REST}
\cvskillline{Outils}{Docker, Git (GitHub / GitLab), CI/CD, PHPStan, PHPUnit, Mercure, Linux}
\cvskillline{Langues}{Français — natif, Anglais — niveau B2}
